\documentclass[11pt,a4paper]{beamer}
\usetheme{Madrid}
\usepackage[utf8]{inputenc}
\usepackage{amsmath}
\usepackage{amsfonts}
\usepackage{amssymb}
\usepackage{graphicx}
\usepackage{xcolor}
\usepackage{tikz}
\usepackage{booktabs}
\usepackage{array}
\usepackage{colortbl}
\usepackage{multicol}

% ============================================================
% EXTREMELY SMALL FONT - MAXIMUM CONTENT PER SLIDE
% ============================================================
\usefonttheme{professionalfonts}
\setbeamerfont{itemize/enumerate body}{size=\tiny}
\setbeamerfont{itemize/enumerate subbody}{size=\tiny}
\setbeamerfont{block body}{size=\tiny}
\setbeamerfont{block title}{size=\scriptsize}
\setbeamerfont{frametitle}{size=\small}
\setbeamerfont{title}{size=\large}
\setbeamerfont{subtitle}{size=\small}
\setbeamerfont{author}{size=\tiny}
\setbeamerfont{date}{size=\tiny}

\renewcommand{\tiny}{\fontsize{5pt}{6pt}\selectfont}
\renewcommand{\scriptsize}{\fontsize{6pt}{7pt}\selectfont}
\renewcommand{\footnotesize}{\fontsize{6.5pt}{8pt}\selectfont}
\renewcommand{\small}{\fontsize{7.5pt}{9pt}\selectfont}
\renewcommand{\normalsize}{\fontsize{8.5pt}{10.5pt}\selectfont}

\newcommand{\examfocus}[1]{\textcolor{orange!80!black}{\textbf{[FOCUS] #1}}}
\newcommand{\trap}[1]{\textcolor{purple}{\textbf{[TRAP] #1}}}
\newcommand{\correct}[1]{\textcolor{green!50!black}{\textbf{[CORRECT] #1}}}
\newcommand{\scenario}[1]{\textcolor{blue!70!black}{\textbf{[SCENARIO] #1}}}
\newcommand{\keyterm}[1]{\textcolor{red!80!black}{\textbf{[KEY] #1}}}
\newcommand{\lawcite}[1]{\textcolor{brown!70!black}{\textbf{[LAW] #1}}}
\newcommand{\redflag}[1]{\textcolor{red!80!black}{\textbf{[RED FLAG: #1]}}}
\newcommand{\bluenote}[1]{\textcolor{blue!80!black}{\textbf{[NOTE: #1]}}}

\title[CAMS Domain C]{CAMS Exam Preparation\\Part C: Building an Anti-Financial Crime\\Compliance Program (30\% of Exam)}
\author{Advanced AML Training Framework}
\date{}

\setbeamercovered{transparent}
\setbeamertemplate{navigation symbols}{}
\setbeamertemplate{frametitle continuation}{}
\setlength{\parskip}{2pt}
\setlength{\itemsep}{1pt}

\begin{document}
	
	% ============================================================
	% TITLE SLIDE
	% ============================================================
	\begin{frame}
		\titlepage
		\centering
		\tiny Domain C: 30 percent of CAMS Exam | Governance, Risk Assessment, CDD, Training, Independent Testing, Investigations
	\end{frame}
	
	% ============================================================
	% SECTION 1: THREE LINES OF DEFENSE (Complete)
	% ============================================================
	
	\begin{frame}{Three Lines of Defense - Complete Framework}
		\begin{block}{Core Concept}
			A foundational governance model ensuring separation between risk-taking, risk oversight, and independent assurance.
		\end{block}
		\begin{tabular}{p{0.18\textwidth}|p{0.32\textwidth}|p{0.45\textwidth}}
			\hline
			\rowcolor{lightgray}
			\textbf{Line} & \textbf{Who} & \textbf{Primary Role} \\
			\hline
			\textbf{First Line} & Business ops, RMs, tellers, loan officers & Own and manage risk; identify anomalies; complete initial CDD; escalate to compliance \\
			\hline
			\textbf{Second Line} & Compliance Dept, BSA Officer, MLRO & Oversee risk; set rules; monitor First Line; investigate; file SARs \\
			\hline
			\textbf{Third Line} & Internal Audit, external auditors & Independent assurance; audit First and Second Lines; report to Board \\
			\hline
		\end{tabular}
		\vspace{0.2cm}
		\redflag{Golden Rule: No single line should both DESIGN and TEST the same controls.}
		\examfocus{Exam tests whether you can identify independence violations across all three lines.}
	\end{frame}
	
	\begin{frame}{First Line - Detailed Responsibilities}
		\begin{block}{First Line (Business Operations) - What They MUST Do}
			The First Line is the front door of AML defense. They interact with customers directly.
		\end{block}
		\begin{multicols}{2}
			\textbf{Specific Duties:}
			\begin{itemize}
				\item Complete initial CIP
				\item Assess expected activity profile
				\item Identify unusual transactions
				\item Escalate concerns to compliance
				\item Respond to compliance RFIs
				\item Maintain customer records
				\item Update changes in customer status
			\end{itemize}
			\columnbreak
			\textbf{Common Red Flags to Spot:}
			\begin{itemize}
				\item Reluctance to provide ID
				\item Unusual transaction patterns
				\item Structuring cash deposits
				\item Changes without explanation
				\item Inconsistent narratives
				\item Third-party transactions
				\item Reluctance to explain source
			\end{itemize}
		\end{multicols}
		\redflag{Critical: First Line cannot close alerts or make SAR decisions. That is Second Line authority.}
	\end{frame}
	
	\begin{frame}{First Line Violation - RM Ignores Red Flags}
		\scenario{RM Tells Analyst to Close Alert Based on Personal Knowledge}
		A relationship manager has known a corporate client for 12 years. The client suddenly makes weekly cash deposits of \$9,500 - just below the \$10,000 CTR threshold. The RM tells compliance: I have known this customer for over a decade. He is trustworthy. Just close the alert.
		
		\vspace{0.1cm}
		\textbf{How to Analyze:}
		\begin{enumerate}
			\item First Line's job: IDENTIFY and ESCALATE - not investigate or close
			\item Second Line's job: INVESTIGATE independently
			\item The deposit pattern (\$9,500 weekly) is structuring
			\item Closing without documentation violates AML program requirements
		\end{enumerate}
		\trap{RM's long relationship is NOT sufficient evidence to close an alert.}
		\correct{Compliance must conduct independent investigation with documented evidence.}
	\end{frame}
	
	\begin{frame}{Second Line - Detailed Responsibilities}
		\begin{block}{Second Line (Compliance) - Core Duties}
			The Second Line must operate independently from the First Line and have authority to override business decisions.
		\end{block}
		\begin{multicols}{2}
			\textbf{Core Duties:}
			\begin{itemize}
				\item Design AML policies
				\item Conduct transaction monitoring
				\item Investigate alerts
				\item Make SAR filing decisions
				\item File SARs with FIU
				\item Provide training
				\item Conduct EDD
				\item Maintain risk assessment
				\item Respond to law enforcement
				\item Report to Board
			\end{itemize}
			\columnbreak
			\textbf{Independence Requirements:}
			\begin{itemize}
				\item Cannot report to business lines
				\item Direct Board access
				\item Override business decisions
				\item Sufficient staffing
				\item No commercial pressure
				\item Access to all information
				\item Independent budget
			\end{itemize}
		\end{multicols}
		\redflag{Commercial importance of a client is NEVER a valid reason to delay or avoid SAR filing.}
	\end{frame}
	
	\begin{frame}{Second Line Failure - No Documentation}
		\scenario{Analyst Closes Alert Without Investigation}
		An analyst reviews an alert for a restaurant customer with 23 cash deposits between \$8,400 and \$9,800 over 45 days. The analyst writes: Restaurant business. Cash is normal. No SAR needed.
		
		\vspace{0.1cm}
		\textbf{How to Analyze:}
		\begin{enumerate}
			\item 23 deposits just below \$10,000 = structuring indicator
			\item Structuring is a federal crime regardless of funds origin
			\item The analyst assumed legitimacy without verification
			\item No documentation = no investigation occurred
		\end{enumerate}
		\trap{Analyst used reasonable business judgment based on industry.}
		\correct{Analysts must obtain supporting documentation and document all investigative steps. Assumptions are not investigations.}
	\end{frame}
	
	\begin{frame}{Third Line - Independence Requirements}
		\begin{block}{Internal Audit - What They CAN and CANNOT Do}
			Audit independence is absolutely critical to program credibility.
		\end{block}
		\begin{tabular}{p{0.45\textwidth}|p{0.45\textwidth}}
			\hline
			\rowcolor{lightgray}
			\textbf{Audit CAN Do} & \textbf{Audit CANNOT Do} \\
			\hline
			Test effectiveness of controls & Design or build controls \\
			Recommend improvements & Operate monitoring systems \\
			Report findings to Board & Make SAR decisions \\
			Review sample of alerts & Approve customer relationships \\
			Assess policy compliance & Perform ongoing monitoring \\
			Evaluate training & Conduct EDD on customers \\
			\hline
		\end{tabular}
		\redflag{Most Common Violation: Audit helping to design scenarios THEN auditing those same scenarios.}
	\end{frame}
	
	\begin{frame}{Third Line Independence Violation}
		\scenario{Audit Designs Controls Then Tests Them}
		Compliance lacks technical expertise. CCO asks Internal Audit for help. Audit spends 3 months co-designing 14 new monitoring scenarios. Six months later, the same audit team tests the system and finds all 14 scenarios effective.
		
		\vspace{0.1cm}
		\textbf{How to Analyze:}
		\begin{enumerate}
			\item Audit's role is to TEST, not DESIGN
			\item Designing controls makes Audit part of Second Line
			\item Auditors cannot objectively test their own work
			\item The audit findings are unreliable
		\end{enumerate}
		\trap{Scenarios were effective, so no violation.}
		\correct{Audit must remain independent. Use external consultants or different internal team with no design involvement.}
	\end{frame}
	
	\begin{frame}{Board of Directors - What They MUST Do}
		\begin{block}{Board Oversight - Ultimate Accountability}
			The Board is ultimately accountable for the AML program. Regulators hold the Board responsible.
		\end{block}
		\begin{multicols}{2}
			\textbf{Board MUST:}
			\begin{itemize}
				\item Approve initial AML program
				\item Approve material changes
				\item Appoint BSA Officer
				\item Receive regular updates
				\item Be informed of deficiencies
				\item Hold management accountable
				\item Review audit findings
				\item Ensure adequate resources
			\end{itemize}
			\columnbreak
			\textbf{Board DOES NOT:}
			\begin{itemize}
				\item File SARs
				\item Review KYC files
				\item Monitor daily alerts
				\item Conduct investigations
				\item Make SAR decisions
				\item Perform testing
				\item Approve transactions
			\end{itemize}
		\end{multicols}
		\redflag{Board approval is required for the AML program. CCO cannot approve their own program.}
	\end{frame}
	
	\begin{frame}{Board Violation - Hiding Deficiencies}
		\scenario{Management Withholds Information from Board}
		Compliance discovers a 14-month monitoring gap affecting \$340 million in wires. Senior management decides not to inform the Board because the issue is operational and will be fixed within 90 days.
		
		\vspace{0.1cm}
		\textbf{How to Analyze:}
		\begin{enumerate}
			\item 14-month gap affecting \$340 million is MATERIAL
			\item Material deficiencies MUST be reported to Board immediately
			\item Management cannot filter or delay reporting
			\item Board needs this information for oversight
		\end{enumerate}
		\trap{Senior management has authority to decide what reaches the Board.}
		\correct{Board must be informed of material deficiencies. Management cannot withhold this information.}
	\end{frame}
	
	\begin{frame}{Board vs. Management - Approval Authority}
		\scenario{CCO Implements Policy Without Board Approval}
		CCO drafts revised AML policy. Legal reviews and approves. CCO implements without presenting to Board. CCO argues: I am the designated BSA Officer with authority to implement policies. Board approval is unnecessary.
		
		\vspace{0.1cm}
		\textbf{How to Analyze:}
		\begin{enumerate}
			\item Board must approve AML program under BSA
			\item Material changes require Board approval
			\item CCO implements, does not have final approval
			\item Policy is not effective without Board approval
		\end{enumerate}
		\trap{CCO has delegated authority to implement policies.}
		\correct{Board must approve AML program and material changes. CCO implements Board-approved program but lacks final approval authority.}
	\end{frame}
	
	\begin{frame}{Senior Management Responsibilities}
		\begin{block}{Senior Management - Execution and Accountability}
			Senior management executes the Board-approved program and is accountable to the Board.
		\end{block}
		\begin{multicols}{2}
			\textbf{Senior Management MUST:}
			\begin{itemize}
				\item Implement Board-approved program
				\item Allocate sufficient resources
				\item Ensure timely SAR filings
				\item Respond to audit findings
				\item Report material issues to Board
				\item Appoint qualified staff
				\item Ensure training effectiveness
			\end{itemize}
			\columnbreak
			\textbf{Senior Management CANNOT:}
			\begin{itemize}
				\item Veto SARs for commercial reasons
				\item Withhold information from Board
				\item Retaliate against whistleblowers
				\item Operate without Board approval
				\item Delegate accountability
				\item Ignore audit findings
			\end{itemize}
		\end{multicols}
		\redflag{Senior management faces personal regulatory risk for AML failures, including fines and imprisonment.}
	\end{frame}
	
	\begin{frame}{AML Program Governance Documentation Requirements}
		\begin{block}{Written Policies, Procedures, and Controls}
			The BSA/AML compliance program must be reduced to writing and approved by the Board.
		\end{block}
		\textbf{Mandatory Governance Documents:}
		\begin{itemize}
			\item Written AML/CFT Policy approved by Board (initial and all material changes)
			\item Written CIP (Customer Identification Program) procedures
			\item Written CDD/EDD procedures (including beneficial ownership)
			\item Written SAR filing procedures (internal decision-making, timeliness)
			\item Written independent testing procedures and scope
			\item Written employee training program (role-based curriculum)
			\item Written risk assessment methodology (Enterprise-Wide Risk Assessment)
			\item Written escalation procedures (compliance to management to Board)
			\item Written recordkeeping procedures (retention periods, secure storage)
		\end{itemize}
		\redflag{Verbal policy or unwritten procedure does NOT satisfy regulatory requirements.}
	\end{frame}
	
	\begin{frame}{Board Meeting Minutes - Documentation Required}
		\scenario{Board Verbally Approves AML Policy Without Minutes}
		Board chair verbally approves updated AML policy during meeting. Secretary does not record vote or discussion in minutes.
		
		\vspace{0.1cm}
		\textbf{How to Analyze:}
		\begin{enumerate}
			\item Unwritten approval is insufficient — approval must be documented in minutes
			\item Minutes must show evidence of meaningful oversight, not rubber-stamping
			\item Discussion of AML risks, resource allocation, and key metrics must appear
			\item Regulators will treat missing minutes as evidence of Board non-engagement
		\end{enumerate}
		\trap{Verbal approval is still approval; minutes are administrative formality.}
		\correct{Board minutes must document AML approval, discussion, and ongoing oversight.}
	\end{frame}
	
	\begin{frame}{CCO/BSA Officer - Independence and Authority}
		\begin{block}{CCO Reporting Line and Authority}
			The CCO/BSA Officer must have sufficient authority, independence, and resources.
		\end{block}
		\textbf{Required Attributes:}
		\begin{itemize}
			\item Direct reporting access to Board or Board Committee (not filtered by business lines)
			\item Authority to allocate resources (staff, technology, training)
			\item Authority to override business decisions regarding customer relationships
			\item No reporting line to revenue-generating functions
			\item Sufficient seniority to influence decisions
			\item Protection from retaliation for exercising duties
		\end{itemize}
		\textbf{Red Flags for Independence Violation:}
		\begin{itemize}
			\item CCO reports to Head of Retail Banking (revenue function)
			\item Business lines control CCO's budget
			\item CCO is removed after filing SAR on high-revenue client
			\item CCO position is vacant for extended period
			\item Same person serves as CCO and Head of Sales
		\end{itemize}
		\redflag{If CCO can be overridden by business lines for commercial reasons, CCO lacks required authority.}
	\end{frame}
	
	\begin{frame}{CCO Reporting Line Violation - Revenue Function}
		\scenario{CCO Reports to Head of International Banking (Revenue)}
		Bank's CCO reporting line is to Head of International Banking, who is compensated based on revenue growth. Head instructs CCO to delay SAR filing on large corporate client because client is negotiating a \$50 million credit facility.
		
		\vspace{0.1cm}
		\textbf{How to Analyze:}
		\begin{enumerate}
			\item CCO must report to Board or independent committee, not revenue-generating management
			\item Revenue pressure on CCO compromises independence
			\item Delaying SAR for commercial reasons is a regulatory violation
			\item Budget reduction following SAR filings suggests retaliation
		\end{enumerate}
		\trap{Reporting to senior management is acceptable as long as CCO has Board access.}
		\correct{CCO must report to independent authority without revenue conflicts.}
	\end{frame}
	
	\begin{frame}{Program Resources - Staffing, Technology, Budget}
		\begin{block}{Adequate Resources = Regulatory Requirement}
			Institutions must provide AML compliance with sufficient resources.
		\end{block}
		\textbf{Resource Areas Requiring Adequacy:}
		\begin{itemize}
			\item Staffing: Sufficient analysts to review all alerts within required timeframes
			\item Technology: Monitoring systems, screening tools, case management, data infrastructure
			\item Training: Ongoing, role-based training for compliance and business staff
			\item External support: Legal, audit, third-party consultants as needed
			\item Budget: Competitive salaries to retain qualified talent
		\end{itemize}
		\textbf{Evidence of Resource Inadequacy:}
		\begin{itemize}
			\item Persistent backlog of uninvestigated alerts (thousands aging >30 days)
			\item High analyst turnover (burnout, underpaid, under-supported)
			\item Manual processes for high-volume tasks that could be automated
			\item Delayed SAR filings due to insufficient reviewers
			\item Board rejects reasonable budget requests from CCO
			\item Staff not qualified for assigned duties
		\end{itemize}
		\redflag{Under-resourcing is NOT a defense. Regulators will cite inadequate resources as a violation.}
	\end{frame}
	
	% ============================================================
	% SECTION 2: ENTERPRISE-WIDE RISK ASSESSMENT (EWRA)
	% ============================================================
	
	\begin{frame}{Enterprise-Wide Risk Assessment - Mandatory Components}
		\begin{block}{Risk Assessment as the Foundation}
			The AML program must be based on a documented, comprehensive risk assessment.
		\end{block}
		\textbf{Required Risk Factors (Must Address All):}
		\begin{multicols}{2}
			\begin{itemize}
				\item Products and services offered
				\item Customer types and segments
				\item Geographic locations (operations and customer base)
				\item Delivery channels (online, branch, mobile, agents)
				\item Transaction volumes and values
				\item Jurisdictions with weak AML regimes
				\item Politically Exposed Persons (PEPs)
				\item Non-profit organization exposure
				\item Correspondent banking relationships
				\item Private banking relationships
				\item Emerging technologies (crypto, etc.)
				\item Previous regulatory findings
			\end{itemize}
		\end{multicols}
		\textbf{Risk Assessment Outputs:}
		\begin{itemize}
			\item Inherent risk score (risk without controls)
			\item Control effectiveness rating
			\item Residual risk score (risk after controls)
			\item Action plan for residual risks exceeding tolerance
		\end{itemize}
		\redflag{EWRA must be UPDATED at least annually or after material changes.}
	\end{frame}
	
	\begin{frame}{Risk Assessment - Inherent vs. Residual Risk}
		\scenario{Bank Has High Inherent Risk But Strong Controls}
		Bank serves high-risk customers (MSBs, NPOs in conflict zones) and operates in high-risk jurisdictions (inherent risk = 8/10). Bank has strong controls: dedicated EDD team, transaction monitoring with low false negatives, monthly Board reporting (control effectiveness = 9/10). Residual risk = 3/10.
		
		\vspace{0.1cm}
		\textbf{How to Analyze:}
		\begin{enumerate}
			\item High inherent risk does NOT automatically mean deficient program
			\item Strong controls can reduce residual risk to acceptable levels
			\item Bank must document both inherent and residual risk
			\item Regulator will test whether controls are as strong as claimed
		\end{enumerate}
		\trap{High inherent risk is always unacceptable regardless of controls.}
		\correct{Risk-based approach accepts higher inherent risk if controls are proportionate and effective.}
	\end{frame}
	
	\begin{frame}{Risk Scoring Models - Numerical Ratings}
		\begin{block}{Quantitative Risk Assessment}
			Many institutions use numerical scores to determine risk categories.
		\end{block}
		\textbf{Typical Risk Scoring Model:}
		\begin{itemize}
			\item Each risk factor (geography, customer type, product/service) scored 1-10
			\item 1-3 = Standard/Low Risk
			\item 4-8 = Medium Risk
			\item 9-10 = High Risk
			\item Composite score = sum or weighted average of all factors
		\end{itemize}
		\textbf{Weighting Considerations:}
		\begin{itemize}
			\item Some factors may be weighted more heavily (e.g., customer type vs. geography)
			\item Combinations can change risk dramatically (e.g., foreign private company + wire transfer capability)
			\item High-risk relationships should not represent too large a segment of portfolio
		\end{itemize}
		\textbf{Review Frequency:}
		\begin{itemize}
			\item At least annually
			\item Before launching new products
			\item After acquisition of another institution
			\item When material changes occur
		\end{itemize}
	\end{frame}
	
	\begin{frame}{Dynamic Customer Risk Reassessment}
		\begin{block}{Risk is Not Static - Must Be Continuously Managed}
			A customer's risk rating may change based on actual activity.
		\end{block}
		\textbf{Triggers for Risk Rating Change:}
		\begin{itemize}
			\item Unusual activity (alerts, cases, SAR filings)
			\item Receipt of law enforcement inquiries (subpoenas)
			\item Transactions that violate economic sanctions
			\item Significant changes in transaction volume or pattern
			\item New information about customer (negative media, change in ownership)
			\item Customer begins using higher-risk products or services
		\end{itemize}
		\textbf{Example:}
		\begin{itemize}
			\item Student checking account = low inherent risk
			\item But if account shows large wires to high-risk jurisdictions = raise risk rating
			\item Conversely, MSB with high inherent risk may have normal activity = lower residual risk
		\end{itemize}
		\redflag{Document ALL risk rating changes with rationale.}
	\end{frame}
	
	% ============================================================
	% SECTION 3: CUSTOMER DUE DILIGENCE (CDD) - ENHANCED
	% ============================================================
	
	\begin{frame}{Customer Due Diligence (CDD) - Four Elements}
		\begin{block}{FATF Recommendation 10 - Core CDD Measures}
			Financial institutions must conduct CDD when establishing business relationships or when suspicion arises.
		\end{block}
		\textbf{Four Required CDD Elements:}
		\begin{enumerate}
			\item Identify the customer and verify identity using reliable, independent source documents
			\item Identify the beneficial owner and take reasonable measures to verify identity
			\item Understand and obtain information on purpose and intended nature of relationship
			\item Conduct ongoing due diligence and scrutiny of transactions to ensure consistency with knowledge of customer, business, and risk profile
		\end{enumerate}
		\textbf{When CDD MUST Be Performed:}
		\begin{itemize}
			\item Establishing business relationships
			\item Carrying out occasional transactions above threshold
			\item Suspicion of money laundering or terrorist financing
			\item Doubts about veracity of previously obtained identification
		\end{itemize}
		\redflag{Anonymous or fictitious name accounts are PROHIBITED.}
	\end{frame}
	
	\begin{frame}{Beneficial Ownership - Verification Requirements}
		\begin{block}{Beneficial Owner Definition}
			The natural person who ultimately owns or controls a customer and/or the natural person on whose behalf a transaction is conducted.
		\end{block}
		\textbf{Verification Requirements for Legal Persons:}
		\begin{itemize}
			\item Identify natural persons with significant ownership (typically 25\% or more)
			\item Identify natural persons with ultimate effective control
			\item Understand ownership and control structure
			\item Use reliable, independent sources (corporate registers, commercial databases)
			\item Do NOT rely solely on written declarations from customer
		\end{itemize}
		\textbf{For Trusts and Similar Arrangements:}
		\begin{itemize}
			\item Identify settlor, trustee, protector (if any), beneficiaries, and any other natural person exercising control
		\end{itemize}
		\textbf{When CDD Cannot Be Performed:}
		\begin{itemize}
			\item Do NOT open account (or close opened accounts)
			\item Consider reporting suspicious activity to authorities
		\end{itemize}
	\end{frame}
	
	\begin{frame}{Enhanced Due Diligence (EDD) - Higher Risk Customers}
		\begin{block}{When Standard CDD Is Not Enough}
			Higher risk circumstances require enhanced measures.
		\end{block}
		\textbf{Customer Risk Factors Requiring EDD:}
		\begin{itemize}
			\item Non-resident customers
			\item Legal persons with nominee shareholders or bearer shares
			\item Cash-intensive businesses
			\item Unusually complex ownership structures
			\item PEPs (foreign and domestic)
		\end{itemize}
		\textbf{Geographic Risk Factors Requiring EDD:}
		\begin{itemize}
			\item Countries with inadequate AML/CFT systems (FATF lists)
			\item Countries subject to sanctions or embargoes
			\item Countries with significant corruption or criminal activity
			\item Countries providing funding/support for terrorism
		\end{itemize}
		\textbf{Product/Service Risk Factors:}
		\begin{itemize}
			\item Private banking
			\item Anonymous transactions (cash, prepaid cards)
			\item Non-face-to-face relationships
			\item Payments from unknown third parties
		\end{itemize}
	\end{frame}
	
	\begin{frame}{EDD - Specific Measures for High-Risk Customers}
		\begin{block}{Enhanced Due Diligence Controls}
			Institutions must apply stronger controls for higher risk relationships.
		\end{block}
		\textbf{Additional Information to Obtain:}
		\begin{itemize}
			\item Source of funds and wealth (documented)
			\item Identifying information on individuals with control (signatories, guarantors)
			\item Financial statements (audited, if available)
			\item Banking references
			\item Description of business operations and expected transaction volume
			\item Explanations for changes in account activity
		\end{itemize}
		\textbf{Required Controls for EDD:}
		\begin{itemize}
			\item Senior management approval to commence or continue relationship
			\item First payment through account in customer's name at bank with similar CDD standards
			\item More frequent monitoring and review of account activity
			\item Periodic re-verification of customer information
		\end{itemize}
	\end{frame}
	
	\begin{frame}{Politically Exposed Persons (PEPs) - Special Requirements}
		\begin{block}{FATF Definition - Two Types of PEPs}
			Individuals entrusted with prominent public functions.
		\end{block}
		\textbf{Foreign PEPs:}
		\begin{itemize}
			\item Heads of state or government
			\item Senior politicians, government, judicial, military officials
			\item Senior executives of state-owned corporations
			\item Important political party officials
			\item Includes family members and close associates
		\end{itemize}
		\textbf{Domestic PEPs:}
		\begin{itemize}
			\item Same categories but within the institution's own country
			\item EU Fourth Directive explicitly covers domestic PEPs
		\end{itemize}
		\textbf{Required Measures for PEPs:}
		\begin{itemize}
			\item Identify PEP status before or during relationship
			\item Senior management approval for relationship
			\item Establish source of wealth and source of funds
			\item Enhanced ongoing monitoring
		\end{itemize}
		\redflag{PEP designation does NOT end when they leave office. Enhanced monitoring typically continues for 12-24 months after position ends.}
	\end{frame}
	
	\begin{frame}{Account Opening and Customer Identification}
		\begin{block}{Natural Persons - Information to Collect}
			Basel Committee General Guide to Account Opening.
		\end{block}
		\textbf{Required Information:}
		\begin{itemize}
			\item Legal name and any other names used (maiden name, alias)
			\item Complete residential address (and business address or PO box based on risk)
			\item Telephone numbers and email address
			\item Date and place of birth, gender
			\item Nationality and residency status
			\item Occupation, position held, name of employer
			\item Official personal identification number
			\item Type of account and nature of relationship
			\item Signature
		\end{itemize}
		\textbf{Verification Methods:}
		\begin{itemize}
			\item Documentary: unexpired official document with photo, official document for date/place of birth
			\item Non-documentary: telephone contact, references from other institutions, public registers, commercial databases
		\end{itemize}
	\end{frame}
	
	\begin{frame}{Legal Persons - Customer Identification}
		\begin{block}{Information to Obtain for Legal Persons/Arrangements}
			Entities require additional due diligence beyond natural persons.
		\end{block}
		\textbf{Required Information:}
		\begin{itemize}
			\item Name, legal form, status, proof of incorporation
			\item Permanent address of principal place of activities
			\item Mailing and registered address
			\item Identity of natural persons authorized to operate account
			\item Official identification number
			\item Powers that regulate and bind the legal person
			\item Identity of beneficial owners
			\item Nature and purpose of activities
			\item Financial situation of entity
			\item Expected use of account (amount, number, type, purpose, frequency)
		\end{itemize}
		\textbf{Verification Methods:}
		\begin{itemize}
			\item Certificate of incorporation, memorandum and articles
			\item Financial statements (audited if available)
			\item Company search, commercial inquiries
			\item Prior bank references
			\item Visit corporate entity where practical
		\end{itemize}
	\end{frame}
	
	\begin{frame}{Consolidated CDD - Enterprise-Wide Approach}
		\begin{block}{Fragmented CDD = Increased Risk}
			A fragmented CDD program significantly increases ML/TF risk.
		\end{block}
		\textbf{Key Principles of Consolidated CDD:}
		\begin{itemize}
			\item Consistent identification and monitoring across business lines globally
			\item Parent-level oversight to capture unusual patterns across subsidiaries
			\item Apply minimum CDD standards to all offices, branches, subsidiaries
			\item When home and host country standards differ, apply HIGHER of the two
		\end{itemize}
		\textbf{What Consolidated CDD Enables:}
		\begin{itemize}
			\item Detect patterns across accounts at different subsidiaries
			\item Identify customer activity that may be low-risk in one region but high-risk in another
			\item Share information between compliance teams across the enterprise
			\item Build complete customer risk profile across all products and geographies
		\end{itemize}
		\redflag{Boundaries or silos in communication should NOT exist within an institution.}
	\end{frame}
	
	% ============================================================
	% SECTION 4: INDEPENDENT TESTING
	% ============================================================
	
	\begin{frame}{Independent Testing - What Must Be Tested}
		\begin{block}{Required Scope - Operational Effectiveness}
			Independent testing is one of the five mandatory BSA pillars.
		\end{block}
		\begin{tabular}{p{0.3\textwidth}|p{0.65\textwidth}}
			\hline
			\rowcolor{lightgray}
			\textbf{Area} & \textbf{Specific Testing Procedures} \\
			\hline
			Transaction Monitoring & Test thresholds using historical SARs; calculate false negative rate; validate alert logic \\
			\hline
			Sanctions Screening & Test fuzzy logic; validate name matching; review override decisions \\
			\hline
			CDD/KYC & Review high-risk samples; verify UBO identification; test EDD documentation \\
			\hline
			SAR Filing & Review timeliness; assess quality; test escalation; validate no tipping-off \\
			\hline
			Training & Assess role-based content; test completion rates; evaluate effectiveness \\
			\hline
			Governance & Review Board approvals; test independence; review escalation documentation \\
			\hline
		\end{tabular}
		\redflag{Not sufficient: Only reviewing policies without testing transactions.}
	\end{frame}
	
	\begin{frame}{Who Can Perform Independent Testing?}
		\begin{block}{Qualification Requirements}
			The tester must meet specific independence and expertise requirements.
		\end{block}
		\begin{tabular}{p{0.45\textwidth}|p{0.45\textwidth}}
			\hline
			\rowcolor{lightgray}
			\textbf{Acceptable} & \textbf{NOT Acceptable} \\
			\hline
			Internal Audit (if independent) & Compliance testing itself \\
			External Audit Firm & Audit that designed controls \\
			Third-Party Consultant & Former compliance staff who designed policies \\
			Independent consultant & Business line employees \\
			No prior involvement & Vendors who built the system \\
			\hline
		\end{tabular}
		\textbf{Requirements:}
		\begin{itemize}
			\item Independence: No involvement in designing or operating controls
			\item Expertise: AML/CFT knowledge and qualifications
			\item Objectivity: No conflicts of interest
			\item Reporting: Reports to Board or Audit Committee, not CCO
		\end{itemize}
	\end{frame}
	
	\begin{frame}{Testing Transaction Monitoring - False Negatives}
		\scenario{Audit Discovers 77 Percent False Negative Rate}
		Auditor tests system by running 500 historical SAR cases through current system. 387 of 500 (77 percent) would NOT generate an alert under current thresholds.
		
		\vspace{0.1cm}
		\textbf{How to Analyze:}
		\begin{enumerate}
			\item 77 percent false negative rate is EXTREMELY high
			\item System would miss over three-quarters of confirmed suspicious activity
			\item Operational burden is NOT a valid excuse
			\item Regulators prioritize detection over efficiency
		\end{enumerate}
		\trap{High alert volume justifies maintaining current thresholds.}
		\correct{System must be effective first. Document deficiency, create remediation plan, recalibrate thresholds.}
	\end{frame}
	
	\begin{frame}{Prior Involvement Trap - Tester Independence}
		\scenario{Tester Previously Worked in Compliance}
		Bank hires external consulting firm for AML audit. Firm assigns individual who was bank's BSA Officer until 8 months ago. This individual personally designed several policies now being audited.
		
		\vspace{0.1cm}
		\textbf{How to Analyze:}
		\begin{enumerate}
			\item Independence is about the PERSON, not just the firm
			\item Person who designed controls cannot audit those controls
			\item 8-month separation does NOT restore objectivity
			\item Firm should assign different personnel
		\end{enumerate}
		\trap{Consultant is external, so independence is automatically satisfied.}
		\correct{Individual tester must have no prior involvement.}
	\end{frame}
	
	\begin{frame}{Audit Reporting Line - CCO Cannot Approve Findings}
		\scenario{Audit Reports to CCO for Approval}
		Audit completes annual AML audit and issues draft report identifying deficiencies. Audit manager provides draft to CCO for approval before finalizing. CCO requests that critical findings be removed. Audit manager agrees.
		
		\vspace{0.1cm}
		\textbf{How to Analyze:}
		\begin{enumerate}
			\item Audit must report findings directly to Board
			\item Management cannot approve or edit audit findings
			\item Removing findings destroys audit independence
			\item Final report is unreliable
		\end{enumerate}
		\trap{CCO is responsible for AML program and should review findings for accuracy.}
		\correct{Audit findings must be reported directly to Board without management filtering.}
	\end{frame}
	
	\begin{frame}{Audit Frequency and Risk-Based Approach}
		\begin{block}{How Often Must Testing Be Performed?}
			Frequency depends on risk profile, but at least annually for most institutions.
		\end{block}
		\textbf{Risk-Based Frequency:}
		\begin{itemize}
			\item High-risk institutions: Semi-annual or quarterly testing recommended
			\item Medium-risk institutions: Annual testing is standard
			\item Low-risk institutions: Annual testing is minimum
		\end{itemize}
		\textbf{Trigger Events for Interim Testing:}
		\begin{itemize}
			\item New high-risk products or services
			\item Significant customer base growth
			\item Merger or acquisition
			\item Regulatory enforcement action
			\item Material system changes
			\item Emergence of new typologies
			\item Known control deficiencies
		\end{itemize}
		\redflag{Annual is minimum, not maximum.}
	\end{frame}
	
	% ============================================================
	% SECTION 5: AML/CFT TRAINING
	% ============================================================
	
	\begin{frame}{AML Training - Role-Based Requirements}
		\begin{block}{Training Must Be Tailored to Specific Roles}
			Different employees need different training based on their responsibilities.
		\end{block}
		\textbf{Role-Based Training Requirements:}
		\begin{itemize}
			\item \textbf{Tellers:} Cash structuring, CTR triggers, identifying suspicious customer behavior, red flags
			\item \textbf{Relationship Managers:} PEP identification, EDD requirements, red flags, source of wealth
			\item \textbf{Board of Directors:} Governance oversight, risk strategy, regulatory expectations, enforcement trends
			\item \textbf{IT/Technology:} Data security, system controls, audit trail requirements, data quality
			\item \textbf{Compliance Staff:} Advanced investigation techniques, SAR writing, regulatory updates, emerging typologies
			\item \textbf{Senior Management:} Program oversight, resource allocation, escalation procedures, regulatory risk
			\item \textbf{New Hires:} Initial training before any customer contact
		\end{itemize}
		\redflag{The same training module delivered to all employees regardless of role is a DEFICIENCY.}
	\end{frame}
	
	\begin{frame}{One-Size-Fits-All Training Failure}
		\scenario{Bank Delivers Same Training to All 5,000 Employees}
		Bank delivers same two-hour online AML training module to all employees: tellers, RMs, Board, IT staff. Training covers technical details of transaction monitoring thresholds.
		
		\vspace{0.1cm}
		\textbf{How to Analyze:}
		\begin{enumerate}
			\item Tellers need training on cash structuring and CTRs, not model calibration
			\item Board members need governance training, not technical threshold details
			\item One-size-fits-all training is ineffective
		\end{enumerate}
		\trap{All employees need same training to ensure consistent institutional coverage.}
		\correct{Training must be tailored to specific roles, responsibilities, and ML/TF risks.}
	\end{frame}
	
	\begin{frame}{Board Training - Governance, Not Operations}
		\scenario{Board Training Covers CTR Filing and Alert Review}
		Bank's annual Board AML training covers: how to identify suspicious transactions in alerts and how to file CTRs.
		
		\vspace{0.1cm}
		\textbf{How to Analyze:}
		\begin{enumerate}
			\item Board members do NOT file CTRs or review alerts (operational)
			\item Board training should cover risk appetite, governance, oversight
			\item Operational training is inappropriate for Board
		\end{enumerate}
		\trap{Board members need to understand all aspects of AML program.}
		\correct{Board training should focus on governance, oversight, risk strategy, and regulatory expectations, not operational tasks.}
	\end{frame}
	
	\begin{frame}{Training Frequency and Documentation}
		\begin{block}{Training Frequency Requirements}
			Initial training AND periodic refresher training are required.
		\end{block}
		\textbf{Frequency Requirements:}
		\begin{itemize}
			\item New hire training: Before any customer contact or within 30 days
			\item Annual refresher training: At least once per calendar year
			\item Role change training: When employee moves to new position
			\item Policy change training: When material changes occur
			\item Remedial training: For employees who fail testing or make errors
		\end{itemize}
		\textbf{Documentation Requirements:}
		\begin{itemize}
			\item Training attendance records (who attended)
			\item Training content (what was covered)
			\item Training dates and duration
			\item Testing results (if applicable)
			\item Certificates of completion
			\item Retention: Minimum 5 years
		\end{itemize}
		\textbf{Training Effectiveness Assessment:}
		\begin{itemize}
			\item Pre- and post-training testing
			\item Scenario-based testing
			\item Observations of employee behavior
			\item Audit of training program
			\item Tracking of repeated errors
		\end{itemize}
		\redflag{Training without effectiveness assessment may not satisfy regulatory requirements.}
	\end{frame}
	
	% ============================================================
	% SECTION 6: KNOW YOUR EMPLOYEE (KYE)
	% ============================================================
	
	\begin{frame}{Know Your Employee (KYE) - Insider Threat Management}
		\begin{block}{Insiders Pose Same Threat as Customers}
			Financial institutions have learned that an insider can pose the same money laundering threat as a customer.
		\end{block}
		\textbf{Key KYE Controls:}
		\begin{itemize}
			\item Background checks on hiring (criminal, credit, employment verification)
			\item Ongoing employee monitoring (not just customer monitoring)
			\item Segregation of duties (no single employee controls end-to-end)
			\item Access controls and audit trails (who accessed what)
			\item Mandatory vacation policies (detect ongoing fraud)
			\item Job rotation for sensitive positions
			\item Employee hotline for reporting suspicious colleagues
			\item Monitoring for lifestyle changes (where permitted by law)
		\end{itemize}
		\redflag{Red flags: Employee living beyond means, refusing vacation, accessing accounts outside job scope, multiple overrides.}
	\end{frame}
	
	\begin{frame}{Employee Background Screening Requirements}
		\begin{block}{Pre-Employment Screening}
			Background screening can prevent hiring of statutorily disqualified individuals.
		\end{block}
		\textbf{What to Screen:}
		\begin{itemize}
			\item Criminal convictions (especially crimes involving dishonesty or money laundering)
			\item Credit history (where permitted)
			\item Employment verification
			\item Education and professional qualifications
			\item References
			\item Bankruptcy searches (for sensitive positions)
			\item ML/TF information sources
		\end{itemize}
		\textbf{Regulatory Prohibitions:}
		\begin{itemize}
			\item Persons convicted of crimes involving dishonesty or money laundering may be prohibited from becoming institution-affiliated parties
			\item Prior written consent of regulator may be required
		\end{itemize}
		\textbf{Ongoing Screening Considerations:}
		\begin{itemize}
			\item Periodic fingerprint checks for sensitive positions
			\item Enhanced screening upon promotion to high-level positions
			\item Monitoring social media (where permitted by law)
		\end{itemize}
	\end{frame}
	
	\begin{frame}{Insider Threat Red Flags and Controls}
		\begin{block}{Detecting Insider Misconduct}
			Employee behavior can indicate complicity in money laundering.
		\end{block}
		\textbf{Employee Red Flags:}
		\begin{itemize}
			\item Employee exaggerates customer credentials in required reports
			\item Employee involved in excessive number of unresolved exceptions
			\item Employee lives lavish lifestyle not supported by salary
			\item Employee frequently overrides internal controls or circumvents policy
			\item Employee uses company resources for private interests
			\item Employee assists transactions where ultimate beneficiary is undisclosed
			\item Employee avoids taking periodic vacations
		\end{itemize}
		\textbf{Mandatory Vacation as a Control:}
		\begin{itemize}
			\item Employee must take consecutive days off (typically 5-10 business days)
			\item Another employee performs duties during absence
			\item Substitute may discover irregularities employee was hiding
			\item Common for wire transfer processors, account opening staff, override approvers
		\end{itemize}
	\end{frame}
	
	% ============================================================
	% SECTION 7: TRANSACTION MONITORING AND CALIBRATION
	% ============================================================
	
	\begin{frame}{False Positives vs. False Negatives}
		\begin{block}{The Core Calibration Trade-Off}
			Every AML monitoring system faces this fundamental decision.
		\end{block}
		\begin{tabular}{p{0.25\textwidth}|p{0.35\textwidth}|p{0.35\textwidth}}
			\hline
			\rowcolor{lightgray}
			\textbf{Term} & \textbf{Definition} & \textbf{Risk} \\
			\hline
			False Positive & Legitimate activity triggers alert & Wasted analyst time; operational inefficiency; analyst fatigue \\
			\hline
			False Negative & Suspicious activity generates no alert & MISSED criminal activity; direct AML deficiency; enforcement risk \\
			\hline
		\end{tabular}
		\vspace{0.2cm}
		\redflag{Which is worse? FALSE NEGATIVES are significantly worse.}
		\bluenote{Regulator Priority: Detection effectiveness > operational efficiency.}
	\end{frame}
	
	\begin{frame}{Over-Filtering Trap - Raise Thresholds Too High}
		\scenario{Director Raises Thresholds by 250 Percent}
		Bank generates 8,000 alerts/month with 96 percent false positives. Director orders all thresholds raised by 250 percent to reduce costs.
		
		\vspace{0.1cm}
		\textbf{How to Analyze:}
		\begin{enumerate}
			\item 250 percent increase is extreme and arbitrary
			\item Many suspicious transactions will fall below thresholds
			\item Director prioritizing cost over effectiveness
		\end{enumerate}
		\trap{Director correctly managing operational efficiency.}
		\correct{Raising thresholds must be data-driven. Validate new thresholds still detect known typologies.}
	\end{frame}
	
	\begin{frame}{Under-Filtering Trap - Too Many Alerts}
		\scenario{System Generates 50,000 Alerts Per Month}
		Bank sets thresholds very low to catch everything. System generates 50,000 alerts/month. Team has 5 analysts, each can review 200 alerts/month. Backlog grows to 40,000 uninvestigated alerts.
		
		\vspace{0.1cm}
		\textbf{How to Analyze:}
		\begin{enumerate}
			\item Team can only review 1,000 alerts/month
			\item 49,000 alerts/month are never reviewed
			\item Unreviewed alerts provide zero detection value
		\end{enumerate}
		\trap{Lower thresholds always better because they catch more activity.}
		\correct{System must be calibrated so team can review all alerts. Unreviewed alerts are a control failure.}
	\end{frame}
	
	\begin{frame}{Model Drift - Performance Degradation}
		\scenario{Model Performance Degrades Over 3.5 Years}
		Bank deployed ML model January 2022 with 72 percent precision and 81 percent recall. June 2025 validation shows 51 percent precision and 62 percent recall. Bank introduced new products, volumes up 180 percent, typologies shifted.
		
		\vspace{0.1cm}
		\textbf{How to Analyze:}
		\begin{enumerate}
			\item Models degrade when behavior/products/typologies change
			\item 3.5 years without recalibration is a significant gap
			\item Annual recalibration or trigger-based retraining required
		\end{enumerate}
		\trap{Model still has recall above 60 percent, which is acceptable.}
		\correct{Bank must establish model monitoring and refresh processes.}
	\end{frame}
	
	% ============================================================
	% SECTION 8: DATA QUALITY AND RECONCILIATION
	% ============================================================
	
	\begin{frame}{Data Quality - The Silent Failure}
		\begin{block}{Core Concept}
			A monitoring system is only as good as the data it receives. Missing data = missing alerts.
		\end{block}
		\textbf{Critical Controls:}
		\begin{itemize}
			\item Reconciliation between source systems and monitoring system (daily)
			\item Data completeness and accuracy checks (automated validation)
			\item Error handling for failed data feeds (no silent failures)
			\item Regular validation that all required transactions are ingested
			\item Data lineage documentation
			\item Data quality metrics and reporting to management
		\end{itemize}
		\redflag{The worst data failure: Silent failure where data is missing but no error messages alert the team.}
	\end{frame}
	
	\begin{frame}{Silent Data Feed Failure - Complete Analysis}
		\scenario{34 Percent of Wires Missing for 14 Months}
		Bank's monitoring system relies on daily feeds from 17 upstream systems. Audit reveals that for 14 months, feed from international wire system silently failed for 34 percent of all wires. Estimated \$340 million in wires not monitored.
		
		\vspace{0.1cm}
		\textbf{How to Analyze:}
		\begin{enumerate}
			\item 34 percent missing data = massive monitoring gap
			\item No alerts for transactions system never received
			\item Absence of reconciliation controls is foundational failure
			\item Bank must remediate AND review missed transactions retrospectively
		\end{enumerate}
		\trap{Misconfiguration was accidental, so no compliance violation.}
		\correct{Bank must have controls to verify complete data ingestion.}
	\end{frame}
	
	\begin{frame}{Data Reconciliation - Required Controls}
		\begin{block}{Reconciliation Controls - What Must Exist}
			Banks must be able to prove that all transactions entered the monitoring system.
		\end{block}
		\textbf{Reconciliation Methodology:}
		\begin{itemize}
			\item Source system transaction count vs. monitoring system intake count (daily)
			\item Source system total dollar volume vs. monitoring system total volume (daily)
			\item Hash total comparisons for data integrity
			\item Exception report for any discrepancies
			\item Timely resolution of exceptions (24-48 hours)
			\item Escalation for unresolved exceptions
		\end{itemize}
		\textbf{Common Data Quality Issues:}
		\begin{itemize}
			\item Missing transactions (data feed failure)
			\item Duplicate transactions (processing errors)
			\item Corrupted data (format issues)
			\item Delayed data (latency beyond acceptable window)
			\item Incomplete data (missing required fields)
			\item Incorrect data (wrong values, mapping errors)
		\end{itemize}
	\end{frame}
	
	% ============================================================
	% SECTION 9: AUTOMATED AML SOLUTIONS
	% ============================================================
	
	\begin{frame}{Automated AML Solutions - Core Capabilities}
		\begin{block}{Technology as an Enabler}
			Manual AML compliance is difficult, if not impossible, for most institutions.
		\end{block}
		\textbf{Key Automated Capabilities:}
		\begin{itemize}
			\item Automated customer verification using third-party databases
			\item Watch list filtering (sanctions, PEPs, adverse media)
			\item Transaction monitoring (scanning and analyzing transactional data)
			\item Automation of regulatory reporting (SARs, CTRs)
			\item Case management (dashboard for KYC, history, investigations, filings)
			\item Audit trail (documenting steps for auditors and supervisors)
		\end{itemize}
		\textbf{Benefits of Automation:}
		\begin{itemize}
			\item Increased efficiency and control
			\item Regulators receive prompt, formatted information
			\item Reflects commitment to meet/exceed compliance requirements
		\end{itemize}
	\end{frame}
	
	\begin{frame}{Selecting an AML System - RFP Process}
		\begin{block}{Request For Proposal (RFP) Method}
			Many institutions use RFPs to select AML software vendors.
		\end{block}
		\textbf{Evaluation Team Composition:}
		\begin{itemize}
			\item Management from compliance, operations, technology, and business departments
			\item Project manager to facilitate
			\item Team reviews and scores all responses
		\end{itemize}
		\textbf{Key Capabilities to Assess:}
		\begin{itemize}
			\item Ability to monitor transactions and identify anomalies
			\item Ability to gather CDD information and score customer responses
			\item Advanced evaluation and analysis of suspicious transactions
			\item View individual alerts within broader customer context
			\item Workflow features (create cases, collaborate, share information)
			\item Automated STR preparation and filing
			\item Standard and ad-hoc reporting
			\item User-friendly updating of risk parameters
			\item Tiered user-rights access
		\end{itemize}
	\end{frame}
	
	\begin{frame}{Case Management Systems - Investigation Tracking}
		\begin{block}{Case Management as Central Hub}
			Effective case management is critical for investigations and audit trails.
		\end{block}
		\textbf{Case Management Features:}
		\begin{itemize}
			\item Dashboard view of all investigations by analyst
			\item Centralized customer KYC and transaction history
			\item Link multiple alerts to single case
			\item Document management for investigation findings
			\item Track SAR/STR filings by customer
			\item Supervisor approval workflow
			\item Audit trail of all case actions
			\item Reporting on investigation volume and outcomes
			\item Aging reports for open cases
		\end{itemize}
		\textbf{Documentation Management:}
		\begin{itemize}
			\item Imaging systems for paperless access to records
			\item Track missing or expired documents
			\item One-stop access to images and compliance documents
		\end{itemize}
	\end{frame}
	
	% ============================================================
	% SECTION 10: NEW PRODUCT RISK ASSESSMENT
	% ============================================================
	
	\begin{frame}{New Products - Pre-Launch Risk Assessment}
		\begin{block}{FATF Recommendation 15 - New Technologies}
			Countries and financial institutions must assess ML/TF risks of new products prior to launch.
		\end{block}
		\textbf{Risk Questions for New Products/Services:}
		\begin{itemize}
			\item Does the product enable significant volumes of transactions rapidly?
			\item Does it allow customers to engage with minimal oversight?
			\item Does it afford significant anonymity to users?
			\item Does it have an especially high transaction or investment value?
			\item Does it allow payments to third parties?
			\item Does it have unusual complexity?
		\end{itemize}
		\textbf{High-Risk Products (Examples):}
		\begin{itemize}
			\item Private banking, offshore international activity
			\item Wire transfer and cash management functions
			\item Prepaid cards, mobile payments, internet-based payment services
			\item Trade financing with unusual pricing features
			\item Payable Through Accounts (PTAs)
		\end{itemize}
		\redflag{Compliance officer must be active participant in new product project teams.}
	\end{frame}
	
	% ============================================================
	% SECTION 11: ALERT ADJUDICATION AND QA
	% ============================================================
	
	\begin{frame}{Alert Adjudication - Investigator Variation Problem}
		\begin{block}{Core Concept}
			Different analysts should reach similar conclusions on similar facts. Wide variation indicates lack of standards.
		\end{block}
		\textbf{Required Controls for Consistency:}
		\begin{itemize}
			\item Written adjudication guidelines for each scenario
			\item Required documentation standards
			\item Quality assurance reviews of closed alerts
			\item Calibration sessions to align analyst judgment
			\item Supervisor approval for certain dispositions
			\item Appeal process for disputed decisions
			\item Regular training on adjudication standards
		\end{itemize}
		\redflag{One analyst closes 95\% of alerts; another escalates 70\% for the same scenario.}
	\end{frame}
	
	\begin{frame}{Quality Assurance Program Requirements}
		\begin{block}{QA Program - Essential Elements}
			Quality assurance is the check on analyst decision-making.
		\end{block}
		\textbf{QA Program Components:}
		\begin{itemize}
			\item Sample size: Sufficient to be statistically valid (typically 5-10\% of alerts)
			\item Sampling methodology: Random plus targeted (high-risk scenarios, specific analysts)
			\item QA reviewer independence: Different from original analyst
			\item QA reviewer qualifications: More experienced than front-line analysts
			\item Dispute resolution process
			\item Remediation process for errors
			\item Trending of QA findings over time
			\item Monthly reporting to compliance management
		\end{itemize}
		\textbf{What QA Reviews:}
		\begin{itemize}
			\item Did analyst obtain required documentation?
			\item Is documentation sufficient to support conclusion?
			\item Was analysis thorough and logical?
			\item Was conclusion reasonable based on evidence?
			\item Were escalation procedures followed?
			\item Was SAR filed when required?
			\item Was tipping-off avoided?
		\end{itemize}
		\redflag{No QA program = no assurance that analysts are making correct decisions.}
	\end{frame}
	
	% ============================================================
	% SECTION 12: SAR FILING AND DECISION-MAKING
	% ============================================================
	
	\begin{frame}{SAR Filing - Timing and Legal Standard}
		\begin{block}{Core Concept}
			SARs must be filed immediately upon suspicion, regardless of transaction value.
		\end{block}
		\textbf{Key Principles:}
		\begin{itemize}
			\item Legal standard = REASONABLE SUSPICION (not certainty, not probable cause)
			\item No minimum monetary threshold for money laundering suspicion
			\item Do NOT delay filing to gather more evidence
			\item File first, continue investigating if needed
			\item Cannot consider commercial impact of filing
			\item Cannot notify customer (tipping-off prohibition)
		\end{itemize}
		\textbf{What Reasonable Suspicion Means:}
		\begin{itemize}
			\item More than a hunch, less than probable cause
			\item Specific articulable facts suggesting possible money laundering
			\item Based on training, experience, and available information
			\item Does not require knowledge of underlying predicate offense
		\end{itemize}
		\redflag{We need more evidence before filing - this is INCORRECT. Reasonable suspicion is enough.}
	\end{frame}
	
	\begin{frame}{STR Decision-Making Process - Factors to Consider}
		\begin{block}{Weighing Aggravating and Mitigating Factors}
			The decision to file an STR results from accumulation of aggravating factors and lack of mitigating factors.
		\end{block}
		\textbf{Aggravating Factors:}
		\begin{itemize}
			\item Structuring transactions below reporting thresholds
			\item Unexplained wire transfers to/from high-risk jurisdictions
			\item Transactions inconsistent with customer's stated business
			\item Customer reluctance to provide information
			\item Negative media or law enforcement inquiries
			\item Multiple alerts on same customer
		\end{itemize}
		\textbf{Mitigating Factors:}
		\begin{itemize}
			\item Documentation supporting legitimate explanation
			\item Consistency with known business activities
			\item Historical pattern of similar activity (not suspicious)
			\item Customer cooperation with RFIs
		\end{itemize}
		\textbf{Documentation Requirement:}
		\begin{itemize}
			\item The decision NOT to file is as important as the decision to file
			\item Document reasoning for both outcomes
		\end{itemize}
	\end{frame}
	
	\begin{frame}{SAR/STR Filing Oversight and Escalation}
		\begin{block}{Oversight of Investigations Process}
			Institutions must have robust policies for oversight of regulatory reporting.
		\end{block}
		\textbf{Escalation Triggers:}
		\begin{itemize}
			\item Customer-facing employee complicity in suspicious activity
			\item Compliance chain-of-command willful blindness
			\item Significant reputational risk to institution
			\item Pattern of filings indicating systemic issue
			\item Potential insider involvement
		\end{itemize}
		\textbf{Reporting to Senior Management/Board:}
		\begin{itemize}
			\item Number of STRs filed (dollar amounts, trends)
			\item Significant compliance deficiencies identified
			\item Law enforcement inquiries related to filings
			\item Emerging typologies observed
		\end{itemize}
		\textbf{Quality Assurance of STRs:}
		\begin{itemize}
			\item STR quality indicates quality of AML program
			\item QA ensures internal consistency
			\item QA reviewers must be qualified and independent
		\end{itemize}
	\end{frame}
	
	\begin{frame}{Account Closure Decision Factors}
		\begin{block}{Closing the Account - Independent Determination}
			Based on internal investigation, institution decides whether to close account.
		\end{block}
		\textbf{Factors to Consider:}
		\begin{itemize}
			\item Legal basis for closing account
			\item Institution's policies (e.g., automatic closure after specified number of STRs)
			\item Seriousness of underlying conduct
			\item Reputational risk to institution if account maintained
			\item Law enforcement requests (to close OR maintain account)
		\end{itemize}
		\textbf{What NOT to Do When Closing:}
		\begin{itemize}
			\item Do NOT tell customer reason if related to SAR (tipping-off)
			\item Do NOT close account to retaliate against whistleblower customer
			\item Do NOT close account without documentation
			\item Do NOT delay returning customer funds
		\end{itemize}
		\textbf{Alternatives to Full Closure:}
		\begin{itemize}
			\item Restrict account activity (no wires, no new deposits)
			\item Enhanced monitoring
			\item Lower transaction limits
			\item Require prior approval for certain transactions
		\end{itemize}
	\end{frame}
	
	\begin{frame}{Communicating with Law Enforcement on STRs}
		\begin{block}{STRs as Financial Intelligence}
			Following STR filing, additional law enforcement notification may be warranted.
		\end{block}
		\textbf{When to Proactively Notify Law Enforcement:}
		\begin{itemize}
			\item Activity involves imminent threat
			\item Activity is time-sensitive
			\item Law enforcement has previously expressed interest in subject
			\item Filing may be lost in high-volume of reports
		\end{itemize}
		\textbf{When Law Enforcement Contacts Institution:}
		\begin{itemize}
			\item May seek underlying information used in investigation
			\item Institution should have policy for responding
			\item Designate specific personnel for LE communication
			\item Document all LE contacts and information provided
		\end{itemize}
		\textbf{Critical:}
		\begin{itemize}
			\item Do NOT disclose STR filing to customer
			\item Do NOT confirm or deny existence of STR
			\item Cooperate as permitted by law
		\end{itemize}
	\end{frame}
	
	% ============================================================
	% SECTION 13: INVESTIGATION SOURCES AND CONDUCT
	% ============================================================
	
	\begin{frame}{Sources of Investigations - Internal Triggers}
		\begin{block}{Proactive and Reactive Investigation Sources}
			Investigations may be initiated from multiple channels.
		\end{block}
		\textbf{Common Investigation Initiators:}
		\begin{itemize}
			\item Regulatory recommendations or official findings
			\item Transaction monitoring rules (automated alerts)
			\item Referrals from customer-facing employees (branch personnel, RMs)
			\item Internal hotlines (ethics, compliance, whistleblower)
			\item Negative media information about customer
			\item Receipt of governmental subpoena or search warrant
		\end{itemize}
		\textbf{Referral Mechanisms:}
		\begin{itemize}
			\item Online forms for branch personnel to report unusual activity
			\item Dedicated AML/CFT compliance email address
			\item Training on what constitutes a referral
		\end{itemize}
		\textbf{Internal Hotlines:}
		\begin{itemize}
		\item	Allow anonymous reporting of employee fraud, code violations, suspicious activity
		\item	Prohibition against retaliation for reporting
		\end{itemize}
	\end{frame}
	
	\begin{frame}{Conducting the Investigation - Key Steps}
		\begin{block}{Effective Financial Investigation Steps}
			Main objective: track movement of money.
		\end{block}
		\textbf{Investigation Steps:}
		\begin{enumerate}
			\item Review internal transactions, customer information, internal documentation
			\item Identify and review external information (internet, public records, media)
			\item Contact business line employees responsible for account relationship
			\item Generate written report documenting relevant findings
		\end{enumerate}
		\textbf{Internal Records to Review:}
		\begin{itemize}
			\item Signature cards, account statements, deposit tickets, checks
			\item Loan records, cashier's checks, traveler's checks, money orders
			\item Wire transfer records, safe deposit box records
			\item Login activity logs, IP addresses (online banking)
		\end{itemize}
		\textbf{Retention Requirements:}
		\begin{itemize}
			\item Typically 5 years for customer account records
		\end{itemize}
	\end{frame}
	
	\begin{frame}{Internet Investigation Tips and Sources}
		\begin{block}{Using the Internet for Investigations}
			External information provides context for internal records.
		\end{block}
		\textbf{Reliable Internet Sources:}
		\begin{itemize}
			\item FATF, Wolfsberg, OECD, OFAC websites
			\item Supervisory authority websites (regulators)
			\item Corporate registrars, electoral rolls
			\item Court lists, judicial decisions, administrative tribunals
			\item Professional oversight bodies (law societies)
		\end{itemize}
		\textbf{Caution with Media Sources:}
		\begin{itemize}
			\item Journalists have biases
			\item In some countries, press is not free
			\item Verify from multiple sources
			\item Consider source's reliability
		\end{itemize}
		\textbf{Search Tips:}
		\begin{itemize}
			\item Prepare plan before searching
			\item Use multiple search engines (no single engine covers entire web)
			\item Use metasearch engines
			\item Use local search engines for foreign countries
			\item Consider commercial databases (fee-based, cross-reference many records)
		\end{itemize}
	\end{frame}
	
	% ============================================================
	% SECTION 14: RECORDKEEPING
	% ============================================================
	
	\begin{frame}{Recordkeeping Requirements - Retention Periods}
		\begin{block}{Mandatory Record Retention}
			Financial institutions must retain specified records for minimum periods.
		\end{block}
		\textbf{Core Retention Requirements:}
		\begin{itemize}
			\item Account opening records (CIP/KYC): 5 years after account closing
			\item Transaction records: 5 years from transaction date
			\item SAR filings and supporting documentation: 5 years from SAR filing date
			\item CTR filings: 5 years from filing date
			\item Funds transfer records (over \$3,000): 5 years
			\item Suspicious activity investigation records (even if no SAR filed): 5 years
			\item Training records: 5 years (track completions, content, dates)
			\item Independent testing reports: 5 years
			\item Board meeting minutes (AML-related): Permanent
		\end{itemize}
		\textbf{What "5 Years" Means:}
		\begin{itemize}
			\item Five years from the date of the last activity or filing
			\item For SARs: five years from SAR filing date
			\item Electronic records must be accessible and readable for full period
		\end{itemize}
		\redflag{Deleting records before retention period expires is a violation.}
	\end{frame}
	
	% ============================================================
	% SECTION 15: REQUESTS FOR INFORMATION (RFIs)
	% ============================================================
	
	\begin{frame}{Requests for Information (RFIs)}
		\begin{block}{Core Concept}
			RFIs are formal requests sent to customers asking for clarification or documentation.
		\end{block}
		\textbf{Best Practices for RFIs:}
		\begin{itemize}
			\item Use when you need more information to resolve an alert
			\item Document the RFI and any response (or lack thereof)
			\item Do NOT disclose that an investigation or SAR is pending (tipping-off)
			\item Failure to respond is itself a red flag
			\item Set reasonable deadline (typically 5-10 business days)
			\item Escalate if no response received
		\end{itemize}
		\redflag{RFIs must be phrased neutrally. Do NOT say "we are investigating you for money laundering."}
	\end{frame}
	
	\begin{frame}{Neutral RFI vs. Tipping-Off Violation}
		\scenario{Analyst Tells Customer About AML Flag}
		Compliance analyst calls customer: Our AML system flagged your deposits as potential structuring. We need to know where you are getting all this cash, or we might have to file a report.
		
		\vspace{0.1cm}
		\textbf{How to Analyze:}
		\begin{enumerate}
			\item Mentioning AML system and potential structuring is tipping-off
			\item Mentioning "file a report" is explicit tipping-off
			\item Customer now knows they are under investigation
		\end{enumerate}
		\trap{Analyst was honest with customer, which is good customer service.}
		\correct{RFIs must be neutral. Example: "We noticed increase in cash deposits. Can you provide documentation explaining source of funds?"}
	\end{frame}
	
	\begin{frame}{Proper RFI Wording and Documentation}
		\begin{block}{Examples of Proper RFI Wording}
			Neutral language that does not tip off the customer.
		\end{block}
		\textbf{Proper Wording Examples:}
		\begin{itemize}
			\item "We noticed an increase in wire transfers from your account. Please provide documentation explaining the business purpose."
			\item "Please confirm the nature of your relationship with [counterparty] and provide any supporting agreements."
			\item "Our records show cash deposits totaling [amount] over [period]. Please provide documentation supporting the source of these funds."
		\end{itemize}
		\textbf{When to Escalate Without RFI:}
		\begin{itemize}
			\item Sending RFI would tip off customer (suspicion is very strong)
			\item Customer has history of destroying records when asked
			\item Law enforcement requests no contact
			\item Immediate SAR filing is more appropriate
		\end{itemize}
	\end{frame}
	
	% ============================================================
	% SECTION 16: TIPPING-OFF AND SAR CONFIDENTIALITY
	% ============================================================
	
	\begin{frame}{Tipping-Off - What Constitutes a Violation}
		\scenario{Teller Tells Customer About Compliance Review}
		Customer asks teller why wire transfer is delayed. Teller says: "Our compliance team is reviewing your account for some unusual activity."
		
		\vspace{0.1cm}
		\textbf{How to Analyze:}
		\begin{enumerate}
			\item Tipping-off includes disclosing that an investigation EXISTS
			\item You do NOT need to mention SAR by name to tip off
			\item Telling customer they are under compliance review IS tipping-off
			\item Violation occurred even though no SAR filed yet
		\end{enumerate}
		\trap{No SAR was filed yet, so no tipping-off occurred.}
		\correct{Tipping-off includes disclosing existence of internal investigation.}
	\end{frame}
	
	\begin{frame}{SAR Confidentiality - Criminal Penalties}
		\begin{block}{Unauthorized Disclosure is a Crime}
			Disclosing SAR information without authorization carries criminal penalties.
		\end{block}
		\textbf{Penalties for Unauthorized Disclosure:}
		\begin{itemize}
			\item Criminal fines up to \$250,000 for individuals
			\item Imprisonment up to 5 years
			\item Civil penalties for institutions
			\item Termination of employment
		\end{itemize}
		\textbf{Who Can Receive SAR Information:}
		\begin{itemize}
			\item FinCEN and other FIUs (by filing)
			\item Law enforcement (via proper channels)
			\item Other opted-in 314(b) institutions (limited)
			\item Legal department (internal need-to-know)
			\item Compliance management (internal need-to-know)
			\item Internal Audit (internal need-to-know)
			\item Board of Directors (oversight)
		\end{itemize}
	\end{frame}
	
	% ============================================================
	% SECTION 17: WHISTLEBLOWER PROTECTIONS
	% ============================================================
	
	\begin{frame}{Whistleblower Protections}
		\begin{block}{Core Concept}
			Laws protect employees who report AML violations from retaliation.
		\end{block}
		\textbf{Key Protections (United States):}
		\begin{itemize}
			\item Bank Secrecy Act (31 U.S.C. Section 5328) prohibits retaliation for reporting BSA violations
			\item Dodd-Frank Act Section 922
			\item FIRREA (financial institution anti-retaliation)
			\item Sarbanes-Oxley Act (public company whistleblowers)
		\end{itemize}
		\textbf{What Constitutes Retaliation:}
		\begin{itemize}
			\item Termination or firing
			\item Demotion or reduction in responsibilities
			\item Suspension without pay
			\item Harassment or hostile work environment
			\item Negative performance reviews (PIPs) without prior issues
			\item Exclusion from meetings or opportunities
		\end{itemize}
		\redflag{Retaliation can result in significant damages and reinstatement.}
	\end{frame}
	
	\begin{frame}{Retaliatory PIP - Complete Analysis}
		\scenario{Analyst Receives PIP After Internal Report}
		Compliance analyst reports internally that senior management is suppressing SAR filings on high-revenue corporate client. Three weeks later, analyst receives performance improvement plan citing communication issues - her first negative review in eight years.
		
		\vspace{0.1cm}
		\textbf{How to Analyze:}
		\begin{enumerate}
			\item Timing (3 weeks after report) suggests retaliation
			\item First negative review in 8 years is suspicious
			\item BSA prohibits retaliation for reporting BSA violations
			\item She can file complaint with DOL/OSHA within 180 days
			\item Internal reports are protected, not just external
		\end{enumerate}
		\trap{Internal reports are not protected; only external reports to FinCEN are protected.}
		\correct{BSA prohibits retaliation for internal reports as well.}
	\end{frame}
	
	% ============================================================
	% SECTION 18: SUPPORTING DEPARTMENTS
	% ============================================================
	
	\begin{frame}{Supporting Departments - Data Privacy and Cybersecurity}
		\begin{block}{AML Compliance Does Not Operate in Isolation}
			Other departments enable effective AFC programs.
		\end{block}
		\textbf{Key Supporting Functions:}
		\begin{itemize}
			\item Data Privacy: Ensures customer data protected while enabling lawful information sharing
			\item Cybersecurity: Protects AML systems from attack; identifies cyber-related red flags
			\item IT/Technology: Maintains transaction monitoring systems, data feeds, audit trails
			\item Human Resources: Employee due diligence, insider threat monitoring
			\item Vendor Management: Third-party due diligence for outsourced AML functions
			\item Legal: Advice on regulatory requirements, subpoena response
			\item Internal Audit: Independent testing of AML controls (Third Line)
			\item Risk Management: Enterprise risk assessment integration
		\end{itemize}
		\redflag{AML cannot function effectively without strong support from these departments.}
	\end{frame}
	
	\begin{frame}{Data Privacy vs. AML - GDPR Exception}
		\scenario{Bank Receives GDPR Right to Erasure Request}
		European bank receives GDPR right to erasure request from customer demanding deletion of all personal data. Customer was subject of SAR filed 18 months ago.
		
		\vspace{0.1cm}
		\textbf{How to Analyze:}
		\begin{enumerate}
			\item GDPR Article 17(3) exempts data retention required by law
			\item AML recordkeeping (minimum 5 years) is legal obligation
			\item Bank can lawfully refuse to delete AML-required records
			\item SAR confidentiality also prohibits acknowledging SAR to customer
		\end{enumerate}
		\trap{GDPR right to erasure is absolute and overrides all other laws.}
		\correct{AML retention requirements are legal exception to GDPR erasure.}
	\end{frame}
	
	\begin{frame}{IT/Technology - System Availability and Integrity}
		\begin{block}{IT's Role in AML Compliance}
			IT maintains the systems that enable AML monitoring. Failures become compliance failures.
		\end{block}
		\textbf{IT Responsibilities for AML:}
		\begin{itemize}
			\item System availability (monitoring must run 24/7/365)
			\item Data integrity (no corruption, accurate processing)
			\item Audit trails (who accessed what, when, from where)
			\item Change management (documented approvals for system changes)
			\item Disaster recovery (backup systems, recovery time objectives)
			\item Security (access controls, preventing unauthorized changes)
			\item Performance (system processes alerts within required timeframes)
		\end{itemize}
		\redflag{IT failures are compliance failures when they impact AML monitoring effectiveness.}
	\end{frame}
	
	% ============================================================
	% SECTION 19: THIRD-PARTY AND VENDOR DUE DILIGENCE
	% ============================================================
	
	\begin{frame}{Third-Party and Vendor Due Diligence}
		\begin{block}{Bank Remains Responsible for Outsourced Functions}
			Vendor risk must be managed. Bank retains ultimate regulatory responsibility.
		\end{block}
		\textbf{Key Requirements:}
		\begin{itemize}
			\item Due diligence on vendors before engagement
			\item Contractual obligations for AML compliance
			\item Right to audit vendor performance
			\item Ongoing monitoring of vendor effectiveness
			\item Bank retains ultimate regulatory responsibility
			\item Cannot outsource accountability
		\end{itemize}
		\textbf{Vendors Requiring AML Due Diligence:}
		\begin{itemize}
			\item Transaction monitoring system providers
			\item Sanctions screening vendors
			\item KYC/CIP verification services
			\item Case management system providers
			\item Training providers
			\item Independent testing firms
		\end{itemize}
		\redflag{Cannot outsource responsibility: Bank is still accountable to regulators.}
	\end{frame}
	
	\begin{frame}{Outsourced Monitoring Failure - Liability}
		\scenario{Vendor System Fails for 8 Months}
		Bank outsources transaction monitoring to third-party vendor. Vendor system fails to generate alerts for 8 months due to technical error. Bank's internal compliance team did not perform any validation of vendor outputs.
		
		\vspace{0.1cm}
		\textbf{How to Analyze:}
		\begin{enumerate}
			\item Bank is ultimately responsible to regulators
			\item Bank cannot delegate liability to vendor
			\item Bank should have performed validation and reconciliation
			\item Bank must remediate (review missed period, file SARs)
		\end{enumerate}
		\trap{Vendor contract makes them liable for monitoring failures.}
		\correct{Bank retains regulatory responsibility regardless of vendor contracts.}
	\end{frame}
	
	\begin{frame}{Vendor Due Diligence Requirements}
		\begin{block}{Pre-Engagement and Ongoing Vendor Diligence}
			Comprehensive due diligence required before and during vendor relationship.
		\end{block}
		\textbf{Pre-Engagement Due Diligence:}
		\begin{itemize}
			\item Vendor background and reputation
			\item Financial stability
			\item Vendor's own AML program
			\item Data security practices
			\item References from other bank clients
			\item Regulatory history (enforcement actions)
			\item Business continuity and disaster recovery
		\end{itemize}
		\textbf{Contract Requirements:}
		\begin{itemize}
			\item Right to audit vendor (on-site and remote)
			\item Service level agreements (uptime, accuracy)
			\item Data ownership and return provisions
			\item Breach notification requirements
			\item Liability and indemnification provisions
		\end{itemize}
	\end{frame}
	
	% ============================================================
	% SECTION 20: ESCALATION AND OFFBOARDING
	% ============================================================
	
	\begin{frame}{Escalation and Offboarding}
		\begin{block}{When Customer Poses Unacceptable Risk}
			Bank must escalate through governance and potentially offboard.
		\end{block}
		\textbf{Escalation Path for High-Risk Customers:}
		\begin{enumerate}
			\item Analyst identifies suspicious activity
			\item Investigation conducted and documented
			\item If SAR filed, notify AML committee
			\item If multiple SARs or high risk, escalate to senior management
			\item Board notification for material risks
			\item Decision to offboard requires senior management approval
		\end{enumerate}
		\textbf{Documentation Required:}
		\begin{itemize}
			\item Summary of suspicious activity
			\item Investigation findings
			\item SARs filed (if any)
			\item Risk assessment of relationship
			\item Attempts to remediate (RFIs, enhanced monitoring)
			\item Recommendation (close or continue)
		\end{itemize}
	\end{frame}
	
	\begin{frame}{Offboarding Decision After Multiple SARs}
		\scenario{Bank Has Filed 7 SARs on Account Over 14 Months}
		Bank filed 7 SARs on account over 14 months. Account shows clear patterns of structured deposits followed by wires to offshore shell companies. RM argues: No conviction has occurred. Client is innocent until proven guilty.
		
		\vspace{0.1cm}
		\textbf{How to Analyze:}
		\begin{enumerate}
			\item Banks are not courts of law
			\item Standard for offboarding is risk, not conviction
			\item 7 SARs on one account is clear evidence of elevated risk
			\item Bank can (and often should) offboard
		\end{enumerate}
		\trap{Customer is innocent until proven guilty, so bank should keep account.}
		\correct{Offboarding is risk management decision, not criminal conviction.}
	\end{frame}
	
	% ============================================================
	% CONCLUSION - DOMAIN C SUMMARY
	% ============================================================
	
	\begin{frame}{Domain C - COMPLETE COVERAGE VERIFICATION}
		\begin{block}{All Topics Now Covered for Domain C (30\% of Exam)}
			The following topics have been covered in this presentation:
		\end{block}
		\begin{multicols}{2}
			\begin{itemize}
				\item Three Lines of Defense
				\item Board governance and minutes
				\item CCO independence and authority
				\item Adequate resources requirement
				\item Enterprise-Wide Risk Assessment (EWRA)
				\item Risk scoring models
				\item Dynamic customer risk reassessment
				\item Customer Due Diligence (CDD) - 4 elements
				\item Beneficial ownership verification
				\item Consolidated CDD
				\item Enhanced Due Diligence (EDD)
				\item PEP requirements
				\item Account opening procedures
				\item Independent testing requirements
				\item Training (role-based)
				\item Know Your Employee (KYE)
				\item Insider threat red flags
				\item Transaction monitoring calibration
				\item Data quality and reconciliation
				\item Automated AML solutions
				\item New product risk assessment
				\item Alert adjudication and QA
				\item SAR filing standard (reasonable suspicion)
				\item STR decision-making process
				\item Account closure decisions
				\item Law enforcement communication
				\item Investigation sources and conduct
				\item Internet investigation tips
				\item Recordkeeping requirements
				\item RFIs (neutral wording)
				\item Tipping-off prohibition
				\item Whistleblower protections
				\item Supporting departments (IT, privacy, cybersecurity)
				\item Vendor due diligence
				\item Escalation and offboarding
			\end{itemize}
		\end{multicols}
		\bluenote{This completes coverage of all Domain C topics for the CAMS exam.}
	\end{frame}
	
	\begin{frame}{Domain C Summary - Top 30 Most Tested Concepts}
		\begin{multicols}{2}
			\begin{enumerate}
				\item Three Lines: First=Business, Second=Compliance, Third=Audit
				\item Independence: Audit cannot design controls it later tests
				\item Board approves program, CCO implements
				\item False negatives worse than false positives
				\item Risk assessment = foundation of program
				\item CDD has four elements (identify, verify, understand, monitor)
				\item Beneficial owner = natural person with control
				\item EDD required for PEPs and high-risk customers
				\item Training must be role-based, not one-size-fits-all
				\item Independent testing must be truly independent
				\item Silent data failure = monitoring gap
				\item Alert adjudication must be documented
				\item SAR standard = reasonable suspicion
				\item Tipping-off includes disclosing investigation exists
				\item RFIs must be neutral phrasing
				\item Whistleblower protections exist for internal reports
				\item Vendor oversight: bank retains responsibility
				\item Offboarding is risk decision, not conviction
				\item Insider threats: monitor employees
				\item Mandatory vacation detects fraud
				\item RMs cannot block SARs for commercial reasons
				\item GDPR has AML retention exception
				\item Data reconciliation prevents silent failures
				\item QA reviews ensure analyst consistency
				\item Senior management faces personal liability
				\item Board must be informed of material deficiencies
				\item Training without effectiveness assessment is insufficient
				\item New products require pre-launch risk assessment
				\item Record retention = 5 years minimum
				\item Under-resourcing is NOT a defense
			\end{enumerate}
		\end{multicols}
	\end{frame}
	
	\begin{frame}{Exam Day Reminders for Domain C}
		\begin{itemize}
			\item Domain C is 30 percent of exam - second largest domain
			\item Three Lines of Defense is foundational - know who does what
			\item Independence is critical - audit cannot audit its own work
			\item SAR standard = reasonable suspicion, not proof
			\item Tipping-off includes disclosing investigation exists
			\item False negatives are worse than false positives
			\item Training must be role-based
			\item Whistleblower protections exist for internal and external reports
			\item Banks cannot outsource regulatory responsibility to vendors
			\item Offboarding is risk decision, not conviction
			\item RMs cannot override compliance SAR decisions
			\item Board must approve AML program
			\item Material deficiencies must be reported to Board
			\item RFIs must be neutral - cannot mention AML or filing
			\item Data reconciliation prevents silent feed failures
			\item QA ensures analyst consistency
			\item Mandatory vacation detects insider fraud
			\item CCO cannot approve their own program - Board approval required
		\end{itemize}
	\end{frame}
	
\end{document}